\documentclass[11pt]{article}

% Change "review" to "final" to generate the final (sometimes called camera-ready) version.
% Change to "preprint" to generate a non-anonymous version with page numbers.
\usepackage[]{acl}
\usepackage{times}
\usepackage{latexsym}

% For proper rendering and hyphenation of words containing Latin characters (including in bib files)
\usepackage[T1]{fontenc}
% For Vietnamese characters
% \usepackage[T5]{fontenc}
% See https://www.latex-project.org/help/documentation/encguide.pdf for other character sets

% This assumes your files are encoded as UTF8
\usepackage[utf8]{inputenc}
\usepackage{microtype}
\usepackage{inconsolata}
\usepackage{graphicx}
\usepackage{array}

% If the title and author information does not fit in the area allocated, uncomment the following
%
%\setlength\titlebox{<dim>}
%
% and set <dim> to something 5cm or larger.

\title{Group X Progress Report:\\My Group's Project Name}


\author{First Author, Second Author, Third Author \\
  \texttt{\{macid1,macid2,macid3\}@mcmaster.ca} }

%\author{
%  \textbf{First Author\textsuperscript{1}},
%  \textbf{Second Author\textsuperscript{1,2}},
%  \textbf{Third T. Author\textsuperscript{1}},
%  \textbf{Fourth Author\textsuperscript{1}},
%\\
%  \textbf{Fifth Author\textsuperscript{1,2}},
%  \textbf{Sixth Author\textsuperscript{1}},
%  \textbf{Seventh Author\textsuperscript{1}},
%  \textbf{Eighth Author \textsuperscript{1,2,3,4}},
%\\
%  \textbf{Ninth Author\textsuperscript{1}},
%  \textbf{Tenth Author\textsuperscript{1}},
%  \textbf{Eleventh E. Author\textsuperscript{1,2,3,4,5}},
%  \textbf{Twelfth Author\textsuperscript{1}},
%\\
%  \textbf{Thirteenth Author\textsuperscript{3}},
%  \textbf{Fourteenth F. Author\textsuperscript{2,4}},
%  \textbf{Fifteenth Author\textsuperscript{1}},
%  \textbf{Sixteenth Author\textsuperscript{1}},
%\\
%  \textbf{Seventeenth S. Author\textsuperscript{4,5}},
%  \textbf{Eighteenth Author\textsuperscript{3,4}},
%  \textbf{Nineteenth N. Author\textsuperscript{2,5}},
%  \textbf{Twentieth Author\textsuperscript{1}}
%\\
%\\
%  \textsuperscript{1}Affiliation 1,
%  \textsuperscript{2}Affiliation 2,
%  \textsuperscript{3}Affiliation 3,
%  \textsuperscript{4}Affiliation 4,
%  \textsuperscript{5}Affiliation 5
%\\
%  \small{
%    \textbf{Correspondence:} \href{mailto:email@domain}{email@domain}
%  }
%}

\begin{document}
\maketitle
% \begin{abstract}
% \end{abstract}

\section{Introduction}

Here, write a brief introduction to the problem you are solving. This can be adapted from your problem description and motivation from the original proposal. This should be around 0.25-0.5 pages.

\section{Related Work}

Here, talk about the related work you encountered for your approach. Cite at least 5 references. Refer to item 2. No one has done exactly your task? Write about the most similar thing you can find. This should be around 0.25-0.5 pages.

\section{Dataset}\label{sec:dataset}

\subsection{Source}

We use the ``Predict Students' Dropout and Academic Success'' dataset from the UCI Machine Learning Repository \citep{realinho2021uci}, which was donated by the Instituto Polit\'{e}cnico de Portalegre in Portugal \citep{realinho2022data}.
It covers undergraduates in 17 degree programs.
The records were compiled from the institution's academic and admissions databases and combined with national macroeconomic indicators.
The donors had already cleaned the data and integer-encoded the categorical fields, so we did not need to annotate anything.

\subsection{Properties}

The dataset has 4,424 instances, one per student.
Each has 36 input features and one target label.
There are no missing values and no duplicate rows.
Table~\ref{tab:feature-groups} groups the features.
Eight are binary and nine are nominal categorical, stored as integer codes with no inherent order.
The rest are numeric or counts.
Semester grades use the Portuguese 0--20 scale.
Admission and previous-qualification grades use a 0--200 scale (observed range 95--190).

\begin{table}[t]
  \centering
  \small
  \setlength{\tabcolsep}{4pt}
  \begin{tabular}{l>{\raggedright\arraybackslash}p{4.3cm}r}
    \hline
    \textbf{Group} & \textbf{Examples} & \textbf{\#} \\
    \hline
    Demographic   & gender, age at enrollment, nationality, displaced & 7 \\
    Socioeconomic & parents' qualification/occupation, debtor, scholarship & 7 \\
    Enrollment    & course, application mode/order, admission grade & 7 \\
    Academic      & per semester: units credited, enrolled, evaluated and approved; grade; units without evaluations & 12 \\
    Macroeconomic & unemployment rate, inflation rate, GDP & 3 \\
    \hline
  \end{tabular}
  \caption{\label{tab:feature-groups}
    The 36 input features, grouped by type.
  }
\end{table}

The target is one of three outcomes: Graduate (2,209; 49.9\%), Dropout (1,421; 32.1\%) or Enrolled (794; 17.9\%).
The classes are imbalanced, so a majority-vote classifier already reaches about 50\% accuracy.
We therefore focus on per-class metrics, as described in Section~\ref{sec:results}.

\subsection{Notable Characteristics}

Following our TA's suggestion to look closely at anomalies, we examined each feature's distribution and found several properties that affect modelling.

\begin{itemize}
  \item \textbf{Zero grades:} 718 students have a first-semester grade of 0, and 870 have a second-semester grade of 0.
    Every one of them also has zero approved units.
    A 0 therefore means ``no units passed or evaluated'', not a real score of 0/20.
    Of this group, 79\% of students are classified as Dropout, so the feature is highly predictive, but its zero value means something different from the rest of the grade scale.
  \item \textbf{Skewed categoricals:} 97.5\% of students share the same nationality, and 89\% are single.
    Only 51 students have educational special needs.
    Several coded features also have many values (father's occupation has 46 codes, mother's qualification has 29), and many of those codes are rare.
  \item \textbf{Age:} age is right-skewed.
    The median is 20 and the range is 17--70, with 662 students older than 30.
  \item \textbf{Macroeconomic features:} each has only about 10 distinct values because they are tied to the year of enrollment.
    They act more like a cohort indicator than a continuous measurement.
\end{itemize}

\subsection{Preprocessing}

The raw CSV needed some cleanup before use.
We removed 73 empty trailing columns left over from the export, stripped a UTF-8 byte-order mark from the first header, and removed a stray tab character from the ``Daytime/evening attendance'' column name.
We then label-encoded the target (Dropout\,=\,0, Enrolled\,=\,1, Graduate\,=\,2).

We kept the zero-grade students and the older students.
Both are real records rather than errors, and the zero-grade group carries a strong dropout signal.
No rows were removed.
We split the data 80/20 into training and test sets, stratified by class.
Any scaling or encoding was fit on the training set only, to avoid leaking test data.
Section~\ref{sec:features} describes how the categorical and numeric features are represented, and Section~\ref{sec:results} describes the evaluation protocol.

The data file (\texttt{data.csv}, about 0.8\,MB) is included in our submission.

\section{Features}\label{sec:features}

Describe any features you used for your model, or how your data was input to your model. Are you doing feature engineering or feature selection? Are you learning embeddings? Is it all part of one neural network? Refer to item 2. This may range from 0.25 pages to 0.5 pages.

\section{Implementation}

Describe your model and implementation here. Refer to item 4. This may take around a page.

\section{Results and Evaluation}\label{sec:results}

How are you evaluating your model? What results do you have so far? What are your baselines? Refer to item 5. This may take around 0.5 pages.

\section{Feedback and Plans}

Write about your plans for the remainder of the project. This should include a discussion of the feedback you received from your TA, and how you plan to improve your approach. Reflect on your implementation and areas for improvement. Refer to item 6. This may be around 0.5 pages.

\section{Template Notes}

You can remove this section or comment it out, as it only contains instructions for how to use this template. You may use subsections in your document as you find appropriate.

\subsection{Tables and figures}

See Table~\ref{citation-guide} for an example of a table and its caption.
See Figure~\ref{fig:experiments} for an example of a figure and its caption.


\begin{figure}[t]
  \includegraphics[width=\columnwidth]{example-image-golden}
  \caption{A figure with a caption that runs for more than one line.
    Example image is usually available through the \texttt{mwe} package
    without even mentioning it in the preamble.}
  \label{fig:experiments}
\end{figure}

\begin{figure*}[t]
  \includegraphics[width=0.48\linewidth]{example-image-a} \hfill
  \includegraphics[width=0.48\linewidth]{example-image-b}
  \caption {A minimal working example to demonstrate how to place
    two images side-by-side.}
\end{figure*}


\subsection{Citations}

\begin{table*}
  \centering
  \begin{tabular}{lll}
    \hline
    \textbf{Output}           & \textbf{natbib command} & \textbf{ACL only command} \\
    \hline
    \citep{Gusfield:97}       & \verb|\citep|           &                           \\
    \citealp{Gusfield:97}     & \verb|\citealp|         &                           \\
    \citet{Gusfield:97}       & \verb|\citet|           &                           \\
    \citeyearpar{Gusfield:97} & \verb|\citeyearpar|     &                           \\
    \citeposs{Gusfield:97}    &                         & \verb|\citeposs|          \\
    \hline
  \end{tabular}
  \caption{\label{citation-guide}
    Citation commands supported by the style file.
  }
\end{table*}

Table~\ref{citation-guide} shows the syntax supported by the style files.
We encourage you to use the natbib styles.
You can use the command \verb|\citet| (cite in text) to get ``author (year)'' citations, like this citation to a paper by \citet{Gusfield:97}.
You can use the command \verb|\citep| (cite in parentheses) to get ``(author, year)'' citations \citep{Gusfield:97}.
You can use the command \verb|\citealp| (alternative cite without parentheses) to get ``author, year'' citations, which is useful for using citations within parentheses (e.g. \citealp{Gusfield:97}).

\subsection{References}

\nocite{Ando2005,andrew2007scalable,rasooli-tetrault-2015}

Many websites where you can find academic papers also allow you to export a bib file for citation or bib formatted entry. Copy this into the \texttt{custom.bib} and you will be able to cite the paper in the \LaTeX{}. You can remove the example entries.

\subsection{Equations}

An example equation is shown below:
\begin{equation}
  \label{eq:example}
  A = \pi r^2
\end{equation}

Labels for equation numbers, sections, subsections, figures and tables
are all defined with the \verb|\label{label}| command and cross references
to them are made with the \verb|\ref{label}| command.
This an example cross-reference to Equation~\ref{eq:example}. You can also write equations inline, like this: $A=\pi r^2$.


% \section*{Limitations}

\section*{Team Contributions}

Write in this section a few sentences describing the contributions of each team member. What did each member work on? Refer to item 7.

% Bibliography entries for the entire Anthology, followed by custom entries
%\bibliography{custom,anthology-overleaf-1,anthology-overleaf-2}

% Custom bibliography entries only
\bibliography{custom}

% \appendix

% \section{Example Appendix}
% \label{sec:appendix}

% This is an appendix.

\end{document}
